\subsection{17.39: no uniform number of system normalizers}
\label{cand:17.39}
\problemstatement{17.39}
For every odd prime \(p\) and \(d\ge1\), there is a supersoluble
metabelian group \(G\) of order \(2p^{2d+1}\), with trivial hypercenter,
for which the least number of system normalizers having trivial
intersection is exactly \(d+1\). Its Frattini subgroup has order \(p\).
Fixing \(p=3\) disproves a uniform bound even at fixed Frattini order.

Write \(V=\F_p^d\), and define
\[
 P=V\times V\times\F_p,\qquad
 (u,v,z)(a,b,c)=(u+a,v+b,z+c+u\cdot b).
\]
These are upper unitriangular block matrices of block sizes \(1,d,1\).
The automorphism \(\alpha(u,v,z)=(u,-v,-z)\) has order two.
Set \(G=P\rtimes\langle t\rangle\), with \(x^t=\alpha(x)\).
The commutator formula is
\[
 [(u,v,z),(a,b,c)]=(0,0,u\cdot b-a\cdot v),
\]
so \(Z(P)=P'=\{(0,0,z)\}=:Z\).
The element \(t\) inverts \(Z\), and every element outside \(P\)
does too. Thus \(Z(G)=1\), and its entire hypercenter is trivial.
The subgroup \(\{(0,v,z)\}\) is abelian normal with abelian quotient,
making \(G\) metabelian.
Coordinate flags in the two copies of \(V\) in \(G/Z\), preceded
by \(Z\), give a normal series of prime-order cyclic factors,
so \(G\) is supersoluble.

A Sylow system consists here of the normal Sylow \(p\)-subgroup
\(P\) and a Sylow \(2\)-subgroup.
Its system normalizer is a conjugate of
\[
 D=N_G(\langle t\rangle)=C_G(t)
   =\{(u,0,0):u\in V\}\times\langle t\rangle.
\]
All such conjugates occur, by Sylow conjugacy, and \(G=DP\)
allows conjugators in \(P\).
Represent an element of \(G\) as \((u,v,z,e)\), meaning
\((u,v,z)t^e\), \(e=0,1\).
For \(g=(a,b,c)\) and \(q=c-a\cdot b\), direct multiplication gives
\[
\begin{split}
 D^g=D_{b,q}
 ={}&\{(u,0,u\cdot b,0):u\in V\}\\
 &{}\cup\{(u,-2b,-2q-u\cdot b,1):u\in V\}.
\end{split}
\]
Every pair \(b,q\) occurs, and different pairs give different subgroups.
Thus there are exactly \(p^{d+1}\) system normalizers.

For \(k\) arbitrary normalizers \(D_{b_i,q_i}\), their intersection
with \(P\) consists of
\[
 \{(u,0,u\cdot b_1,0):
       u\cdot(b_i-b_1)=0,\ 2\le i\le k\}.
\]
At most \(k-1\) independent equations are imposed on \(V\);
for \(k\le d\) this intersection is nontrivial.
Conversely take \(D_{0,0}\) and \(D_{e_i,0}\) for all \(d\) standard
basis vectors. Their common part in \(P\) is trivial.
Already the first two have no common element outside \(P\), whose
second coordinates would have to be both \(0\) and \(-2e_1\).
Their full intersection is therefore trivial, proving the exact minimum.

To locate the example relative to the known Frattini-free positive
case \cite{system-normalizers}, note that \(\Phi(G)=Z\).
Indeed \(P^p=1\), by
\[
 (u,v,z)^m=(mu,mv,mz+\tbinom m2u\cdot v),
\]
so \(\Phi(P)=P'P^p=Z\).
If a maximal subgroup \(M\) of \(G\) missed this normal subgroup,
then \(G=M\Phi(P)\) and
\(P=(P\cap M)\Phi(P)\), contrary to the nongenerator property
of \(\Phi(P)\). Thus \(Z\le\Phi(G)\).
On the other hand
\(G/Z\cong V\times(V\rtimes C_2)\), with inversion on the second
copy. The coordinate-hyperplane subgroups in either copy, extended
by the other factors, and the index-two subgroup \(V\times V\),
are maximal with trivial total intersection.
Their preimages give \(\Phi(G)\le Z\).
The fixed nontrivial Frattini subgroup therefore explains why
the prior positive result does not apply.
See \artifact{research/17.39-proof.md}.
